\begin{scanpage}{0134}{116a}
nonuniform part of the current on each side of the strip also has the form of waves which diverge from edges 1 and 2, and are an ``analytical extension'' of the corresponding terms in Equations (20.02) and (20.03). The current wave encountering the opposite edge is reflected from it. Or else one may say that each of the cylindrical waves propagated from edge 1 or 2 undergoes diffraction by the opposite edge (secondary diffraction).

If the strip's width is sufficiently large in comparison with the wavelength, then one may approximately assume that the oncoming current wave near the strip's edge will be the same as on a corresponding half-plane excited by a linear source, the moment of which is selected in a definite way. It is also obvious that the current waves reflected from the edge will also coincide. Consequently, the problem of secondary diffraction by a strip may be reduced to the problem of the diffraction of a cylindrical wave by a half-plane.

The field created at the point $P$ by a current filament parallel to the half-plane's edge and passing through the point $Q$ (Figure 46) may be found by means of the reciprocity principle. In the case of E-polarization, it is determined by the relationship
\[E_z=\frac{p_z}{p_{0z}}E_z(Q),\tag{20.06}\]
and in the case of H-polarization
\[H_z=\frac{m_z}{m_{0z}}H_z(Q).\tag{20.07}\]
Here $p_z$ $(m_z)$ is the electric (magnetic) moment of the current filament passing through the $Q$; $p_{0z}$ $(m_{0z})$ is the moment of the auxiliary current filament passing through the point $P$ with the coordinates $(\phi'',R)$, and $H_z(Q)$ or $E_z(Q)$ is the field created by the auxiliary filament at the point $Q$.

Now let us remove the auxiliary current filament to such a distance that the cylindrical wave arriving from it may be considered
\end{scanpage}
